\section{Discussion}\label{sec:discussion}

\subsection{Two degradations matter, and the failure differs}
Across two sweeps of the same tests, only blur and low resolution produced large error growth, and only at their most severe levels did they produce a falsely reassuring heart rate. The other four axes, over the ranges tested, either did nothing measurable (JPEG compression, grid contrast, uneven lighting) or caused pictures to be unreadable rather than wrong (tilt). This distinction is practical: a system that returns nothing can ask the user to retake the photograph, whereas a system that returns a confident wrong number cannot. It also means that a quality safeguard need not treat all degradation alike.

The pattern within the interval results suggests a mechanism, which we state as a hypothesis because we did not instrument the digitiser. Blur and downscaling both smooth the trace. The narrowest feature, the QRS complex, was affected most (median error 117.6\,ms at blur level~5), QT and QTc next, and the slow P wave least (PR error did not rise monotonically). That ordering is what one expects if smoothing widens sharp features and blurs their boundaries, and it is consistent with the digitiser reconstructing a thicker trace. Tilt, by contrast, is a global geometric change that the digitiser's layout step may already correct for, and JPEG compression, lighting and grid fading alter appearance without removing the trace's edges at the severities tested.

\subsection{Where a quality gate would have to sit}
Interval error began to rise significantly at level~2 of blur and level~3 of resolution, and then rose steeply; falsely reassuring heart rates appeared only at level~5 of both. In the units of this study, level~4 of resolution is a page reduced to a quarter of its linear size and restored, and level~4 of blur is a 15\,px kernel on a page 2{,}200\,px wide. These are conditions of the simulation, not universal thresholds: a real photograph's blur and resolution depend on the camera, distance and lighting, and an operational gate would have to be calibrated on real photographs. What the results do support is the shape of the requirement: a gate should be specific to sharpness and effective resolution, need not fire on moderate tilt, compression or lighting, and cannot rely on a gradual early warning for heart rate, because under low resolution the heart-rate error is flat and then jumps at level~5.

\subsection{Interval measurements need a stated uncertainty}
Heart rate is measured robustly, with a clean-image bias of 0.26\bpm\ and 95\,\% limits of agreement of $-3.1$ to $+3.6$\bpm. PR, QRS, QT and QTc are not: on undamaged images they disagreed with ground truth by 11--22\,ms on average, with limits of agreement spanning 51--110\,ms. The prolonged-QTc thresholds (450 and 460\,ms) sit inside that band for many patients, which is why a naive count of falsely reassuring QTc classifications was inflated (26 events, 20 of them on records already misclassified on clean images). A patient-facing statement built on such a value should not carry the same confidence as one built on heart rate; it needs either a more precise delineation method or an explicit uncertainty band. This is a concrete argument for uncertainty-qualified explanation, and it is a result about the measurement, independent of the degradation.

The partial rescue from a 500\,Hz ground truth is itself informative. It shows that some of what looks like digitiser error in interval studies is ground-truth error, and that a study using a low-rate reference will overstate the error floor. It also shows the reverse for heart rate: because ground-truth heart rate differs by 0.75\bpm\ between 100 and 500\,Hz, a clean-image heart-rate error of 0.53\bpm\ cannot be interpreted as digitiser error alone.

\subsection{Reproducibility is a result, not just a hygiene issue}
The heart-rate findings came out identically in the second sweep: the same six tests were falsely reassuring, and 99.3\,\% of the tests readable in both agreed to within 0.15\bpm. This is a computational repeat, not an independent replication, because the records, pages and code were the same; it shows that the pipeline is deterministic, not that the effect would recur on new data. What did change was which records could be read: seven records changed status, including all four that failed on their clean image in Sweep~1. A study that reported only one sweep would have concluded either that four records are inherently undigitisable or that the digitiser is highly reliable; the truth is that reliability depended on the tool build. This argues for pinning tool versions, recording them, and reporting exclusions as properties of a build, not of an image.

\subsection{Relation to earlier work}
Digitisation studies commonly report signal-to-noise ratio \citep{ref2,ref5}. Our study, like the observation of \citet{ref2} that this metric does not directly reflect clinical measurements, measures clinical values instead; our own attempt to report SNR failed for a mundane reason (unaligned sampling rates), which is a small reminder of how easily a signal-level metric can be computed wrongly and go unnoticed. Quality-assessment work asks whether a signal is usable \citep{ref12,ref13,ref14}; we measure what happens downstream when it is not. And the abstention literature reminds us that declining to answer is not automatically safe \citep{ref26}. Our results suggest \emph{when} an image-quality gate might need to decline (severe blur or low resolution), but they do not test what a patient does after being asked to retake a photograph.

\subsection{What we would do differently}
(i) Pin and record every tool version from the first run. (ii) Save digitised signals from the start, so that a second, more expensive sweep is not needed for new metrics. (iii) Check the clean-image baseline of every new metric \emph{before} interpreting degradation effects; both the interval noise floor and the invalid SNR would have been seen immediately. (iv) Compute ground truth at the highest available sampling rate. (v) Align signals in time and rate before any sample-wise comparison.
